\documentclass[11pt,spanish]{article}
\renewcommand{\sfdefault}{lmss}
\renewcommand{\familydefault}{\rmdefault}
\usepackage[utf8]{inputenc}
\usepackage{geometry}
\geometry{verbose,tmargin=2cm,bmargin=2cm,lmargin=2cm,rmargin=2cm,columnsep=1cm,marginparwidth=1cm}
\usepackage{fancyhdr}
\pagestyle{fancy}
\setlength{\headheight}{13.6pt}
\usepackage[skip=\bigskipamount]{parskip}
\usepackage{array}
\usepackage{microtype}
\usepackage{hyperref}
\hypersetup{hidelinks}
\usepackage{pifont}
\usepackage{titling}
\setlength{\droptitle}{-2cm}
\usepackage{xcolor}
\usepackage{enumitem}
\usepackage[most]{tcolorbox}
\usepackage{amssymb}

\definecolor{darkpastelblue}{rgb}{0.47, 0.62, 0.8}
\definecolor{amber}{rgb}{1.0, 0.75, 0.0}
\definecolor{rojo7a3c3d}{rgb}{0.478, 0.235, 0.239}
\definecolor{verde587246}{rgb}{0.345, 0.447, 0.275}
\definecolor{antiquefuchsia}{rgb}{0.57, 0.36, 0.51}
\definecolor{pastelgreen}{rgb}{0.82,0.92,0.78}
\definecolor{indiagreen}{rgb}{0.07,0.53,0.03}

\newcommand{\blank}[1][4cm]{\rule{#1}{0.4pt}}
\newcommand{\cuadro}{\(\square\)}
\newcommand{\cf}[2]{\colorbox{#1}{\textcolor{white}{\strut #2}}}
\ifrespuestas
\newcommand{\espaciosalida}{0.10cm}
\else
\newcommand{\espaciosalida}{0.65cm}
\fi

\newtcolorbox{ejercicio}[1]{
colback=antiquefuchsia!5!white,
colframe=antiquefuchsia!85!black,
fonttitle=\bfseries,
title={#1},
breakable}

\newtcolorbox{observacion}[1]{
colback=amber!5!white,
colframe=amber!85!black,
fonttitle=\bfseries,
title={#1},
breakable}

\newtcolorbox{respuesta}[1]{
colback=pastelgreen!35!white,
colframe=indiagreen!85!black,
fonttitle=\bfseries,
title={#1},
breakable}

\setlength{\parskip}{\bigskipamount}
\setlength{\parindent}{0pt}

\AtBeginDocument{
  \def\labelitemi{\(\star\)}
  \def\labelitemii{\(\ast\)}
  \def\labelitemiii{\(\bullet\)}
  \def\labelitemiv{\(\circ\)}
}

\usepackage{babel}
\addto\shorthandsspanish{\spanishdeactivate{~<>.}}
\deactivatequoting

\begin{document}

\lhead{PDS}
\rhead{Cuadernillo 02}
\title{Programación de Sistemas\\
Unidades 1 y 2: Estructura del compilador y análisis léxico\\ \bigskip
\cf{rojo7a3c3d}{\textbf{Cuadernillo de ejercicios 02}}\\
\bigskip
\Large\textbf{Del software de sistemas al análisis léxico}
\ifrespuestas\\[0.4em]\large Versión para docente\fi}
\author{Manuel Soto Romero \and Araceli Liliana Reyes Cabello}
\date{Semestre 2026-2 \\ Colegio de Ciencia y Tecnología UACM}

\maketitle
\renewcommand*{\thefootnote}{\arabic{footnote}}
\setcounter{footnote}{0}

Este cuadernillo acompaña las notas 02 y 03. Su propósito es practicar la estructura general de un compilador, la función y comunicación de sus fases y los fundamentos del análisis léxico. Las actividades avanzan desde identificar entradas y salidas hasta seguir la transformación de caracteres en tokens y distinguir token, lexema, patrón y atributo.

\begin{observacion}{Observación 1. Límite de trabajo}
Responde únicamente con lo estudiado en las notas 02 y 03. Los fragmentos de código son ilustrativos: no definen la sintaxis de \textsc{IMP}. No necesitas escribir expresiones regulares, construir autómatas, proponer reglas oficiales de \textsc{IMP} ni diseñar estrategias de recuperación de errores.
\end{observacion}

\begin{observacion}{Observación 2. Ruta sugerida}
Los ejercicios \textbf{1, 3, 5, 7, 8 y 9 son esenciales}: cubren los contratos del compilador y la transformación de caracteres en tokens. Los ejercicios \textbf{2, 4, 6 y 10 son de extensión}: profundizan o integran lo aprendido. Si el tiempo es limitado, completa primero los esenciales; ningún ejercicio se elimina del cuadernillo.
\end{observacion}

\section*{\cf{verde587246}{1. Estructura y fases del compilador}}

\begin{ejercicio}{Ejercicio 1. Reconocer los contratos entre fases \hfill {\small [Esencial]}}
Completa la columna de salida con las expresiones del banco. Cada expresión se usa una sola vez.

\begin{center}
\small
\begin{tabular}{|c|c|c|}
\hline
programa destino & representación estructurada & secuencia de tokens \tabularnewline
\hline
representación intermedia equivalente & representación validada & representación intermedia \tabularnewline
\hline
\end{tabular}
\end{center}

\begin{center}
\small
\renewcommand{\arraystretch}{1.35}
\begin{tabular}{|p{0.20\textwidth}|p{0.30\textwidth}|p{0.35\textwidth}|}
\hline
\textbf{Fase} & \textbf{Entrada principal} & \textbf{Salida} \tabularnewline
\hline
Análisis léxico & Caracteres del programa fuente & \tabularnewline[\espaciosalida]
\hline
Análisis sintáctico & Secuencia de tokens & \tabularnewline[\espaciosalida]
\hline
Análisis semántico & Representación estructurada e información necesaria & \tabularnewline[\espaciosalida]
\hline
Generación de representación intermedia & Representación validada & \tabularnewline[\espaciosalida]
\hline
Optimización & Representación intermedia & \tabularnewline[\espaciosalida]
\hline
Generación de código & Representación validada o intermedia & \tabularnewline[\espaciosalida]
\hline
\end{tabular}
\end{center}
\end{ejercicio}

\ifrespuestas
\begin{respuesta}{Respuesta 1}
En orden: secuencia de tokens; representación estructurada; representación validada; representación intermedia; representación intermedia equivalente; programa destino.
\end{respuesta}
\fi

\ifrespuestas\else\newpage\fi
\begin{ejercicio}{Ejercicio 2. Análisis o síntesis \hfill {\small [Extensión]}}
Escribe \textbf{A} si la tarea corresponde a la etapa general de análisis y \textbf{S} si corresponde a la etapa general de síntesis.

\begin{center}
\small
\begin{tabular}{|p{0.76\textwidth}|p{0.10\textwidth}|}
\hline
\textbf{Tarea} & \textbf{A/S} \tabularnewline
\hline
Agrupar caracteres para producir tokens. & \tabularnewline[0.55cm]
\hline
Reconocer la estructura de una secuencia de tokens. & \tabularnewline[0.55cm]
\hline
Comprobar restricciones de declaración y tipos. & \tabularnewline[0.55cm]
\hline
Construir una representación intermedia a partir de una representación comprobada. & \tabularnewline[0.55cm]
\hline
Transformar una representación para mejorar una característica sin cambiar el comportamiento definido. & \tabularnewline[0.55cm]
\hline
Producir instrucciones o construcciones del lenguaje destino. & \tabularnewline[0.55cm]
\hline
\end{tabular}
\end{center}

Explica por qué la representación producida por el análisis es necesaria para iniciar la síntesis:

\ifrespuestas\else
\blank[14cm]

\blank[14cm]
\fi
\end{ejercicio}

\ifrespuestas
\begin{respuesta}{Respuesta 2}
Orden esperado: A, A, A, S, S, S.

La representación producida por el análisis comunica una versión estructurada y comprobada del programa. La síntesis utiliza esa información para transformar el programa sin tener que volver a interpretar directamente el flujo original de caracteres.
\end{respuesta}
\fi

\section*{\cf{verde587246}{2. Responsabilidades y diagnósticos}}

\begin{ejercicio}{Ejercicio 3. ¿Qué fase debe responder? \hfill {\small [Esencial]}}
Escribe \textbf{L} para análisis léxico, \textbf{S} para análisis sintáctico o \textbf{M} para análisis semántico.

\begin{center}
\small
\begin{tabular}{|p{0.77\textwidth}|p{0.09\textwidth}|}
\hline
\textbf{Situación} & \textbf{L/S/M} \tabularnewline
\hline
Ningún patrón de token permite reconocer el inicio de los caracteres restantes. & \tabularnewline[0.62cm]
\hline
Los tokens fueron reconocidos, pero su secuencia no puede organizarse de acuerdo con la gramática. & \tabularnewline[0.62cm]
\hline
Un identificador aparece en una estructura válida, pero no satisface una regla de declaración. & \tabularnewline[0.62cm]
\hline
Dos operandos forman una expresión sintácticamente válida, pero sus tipos no son compatibles con la operación. & \tabularnewline[0.62cm]
\hline
Una secuencia de caracteres se reconoce como un identificador. & \tabularnewline[0.62cm]
\hline
Una secuencia de tokens se reconoce como una asignación. & \tabularnewline[0.62cm]
\hline
\end{tabular}
\end{center}

¿Por qué el analizador léxico no debe decidir por sí solo si un identificador fue declarado correctamente?

\ifrespuestas\else
\blank[14cm]

\blank[14cm]
\fi
\end{ejercicio}

\ifrespuestas
\begin{respuesta}{Respuesta 3}
Orden esperado: L, S, M, M, L, S.

El analizador léxico reconoce y clasifica el lexema como identificador. Comprobar una declaración requiere considerar la estructura del programa y la información semántica disponible; por tanto, corresponde a una fase posterior.
\end{respuesta}
\fi

\begin{ejercicio}{Ejercicio 4. Delimitación de responsabilidades \hfill {\small [Extensión]}}
Marca cada afirmación como \textbf{verdadera} o \textbf{falsa}. Corrige las afirmaciones falsas.

\begin{enumerate}[leftmargin=*]
    \item El analizador sintáctico recibe una secuencia de tokens. \hfill V \cuadro\quad F \cuadro
    \item El analizador semántico comprueba restricciones que no dependen únicamente de la forma gramatical. \hfill V \cuadro\quad F \cuadro
    \item El analizador léxico determina la estructura completa de las expresiones. \hfill V \cuadro\quad F \cuadro
    \item La optimización puede cambiar el resultado del programa si con ello reduce su tamaño. \hfill V \cuadro\quad F \cuadro
    \item La generación de código produce una representación en el lenguaje destino. \hfill V \cuadro\quad F \cuadro
    \item Todos los compiladores deben tener exactamente una representación intermedia y un optimizador. \hfill V \cuadro\quad F \cuadro
\end{enumerate}

\ifrespuestas\else
\textbf{Correcciones:}

\blank[14cm]

\blank[14cm]

\blank[14cm]
\fi
\end{ejercicio}

\ifrespuestas
\begin{respuesta}{Respuesta 4}
1. Verdadera. 2. Verdadera. 3. Falsa. 4. Falsa. 5. Verdadera. 6. Falsa.

Correcciones:
\begin{itemize}[leftmargin=*]
    \item El analizador \textbf{sintáctico} determina cómo se organizan los tokens en estructuras; el léxico reconoce unidades a partir de caracteres.
    \item La optimización debe conservar el comportamiento definido del programa.
    \item Un compilador puede usar varias representaciones intermedias, repetir optimizaciones o no incluir un optimizador en una implementación sencilla.
\end{itemize}
\end{respuesta}
\fi

\section*{\cf{verde587246}{3. Integración del compilador del curso}}

\begin{ejercicio}{Ejercicio 5. Completar el recorrido de \textsc{IMP} a \textsc{C} \hfill {\small [Esencial]}}
Completa el flujo con las cuatro fases del banco.

\begin{center}
\small
\begin{tabular}{|c|c|c|c|}
\hline
generación de código & análisis léxico & verificación semántica & análisis sintáctico \tabularnewline
\hline
\end{tabular}
\end{center}

\begin{center}
\renewcommand{\arraystretch}{1.55}
\begin{tabular}{c@{\(\longrightarrow\)}c@{\(\longrightarrow\)}c}
programa \textsc{IMP} & \blank[3.3cm] & tokens \\[0.3cm]
tokens & \blank[3.3cm] & estructura o ASA\footnote{ASA significa \emph{árbol de sintaxis abstracta}. La denominación proviene del inglés \emph{Abstract Syntax Tree}, abreviado AST.} \\[0.3cm]
estructura o ASA & \blank[3.3cm] & representación validada \\[0.3cm]
representación validada & \blank[3.3cm] & programa \textsc{C}
\end{tabular}
\end{center}

Responde brevemente.

\begin{enumerate}[leftmargin=*]
    \item ¿Qué fases serán apoyadas por Lark? \ifrespuestas\else\blank[7cm]\fi
    \item ¿Qué fases se implementarán explícitamente mediante recorridos en Python? \ifrespuestas\else\blank[6cm]\fi
    \item ¿El programa \textsc{C} producido ya es un archivo ejecutable? Explica.

    \ifrespuestas\else
    \blank[14cm]

    \blank[14cm]
    \fi

    \item ¿Dónde puede ubicarse una representación intermedia en este recorrido?

    \ifrespuestas\else
    \blank[14cm]

    \blank[14cm]
    \fi
\end{enumerate}
\end{ejercicio}

\ifrespuestas
\begin{respuesta}{Respuesta 5}
En orden: análisis léxico; análisis sintáctico; verificación semántica; generación de código.

\begin{enumerate}[leftmargin=*]
    \item Lark apoyará el análisis léxico y el análisis sintáctico.
    \item La verificación semántica y la generación de código se implementarán explícitamente en Python mediante recorridos sobre la representación construida.
    \item No. El destino de nuestro compilador es un programa en \textsc{C}; después debe intervenir un compilador de \textsc{C} para continuar la cadena hacia un ejecutable.
    \item Puede ubicarse después de obtener una representación validada y antes de la generación de código. Su forma concreta se definirá en otra nota.
\end{enumerate}
\end{respuesta}
\fi

\begin{ejercicio}{Ejercicio 6. Interfaces y cambios \hfill {\small [Extensión]}}
Un equipo define esta interfaz:

\begin{center}
\fbox{analizador léxico}
\(\longrightarrow\)
\fbox{secuencia de tokens}
\(\longrightarrow\)
\fbox{analizador sintáctico}
\end{center}

Responde.

\begin{enumerate}[leftmargin=*]
    \item ¿Qué debe conocer el analizador sintáctico sobre la salida del léxico?

    \ifrespuestas\else
    \blank[14cm]

    \blank[14cm]
    \fi

    \item Si cambia la forma de representar los tokens, ¿qué componente debe revisarse además del analizador léxico y por qué?

    \ifrespuestas\else
    \blank[14cm]

    \blank[14cm]
    \fi

    \item Explica por qué una interfaz clara ayuda a probar el compilador por fases.

    \ifrespuestas\else
    \blank[14cm]

    \blank[14cm]
    \fi
\end{enumerate}
\end{ejercicio}

\ifrespuestas
\begin{respuesta}{Respuesta 6}
\begin{enumerate}[leftmargin=*]
    \item Debe conocer el formato de cada token: el nombre que representa la categoría y, cuando se requiera, el atributo comunicado.
    \item Debe revisarse el analizador sintáctico porque consume esa representación. También deben revisarse otros componentes que utilicen directamente los atributos modificados.
    \item Una interfaz clara permite preparar entradas controladas, observar salidas y verificar cada responsabilidad sin ejecutar necesariamente todo el compilador.
\end{enumerate}
\end{respuesta}
\fi

\section*{\cf{verde587246}{4. Fundamentos del análisis léxico}}

\begin{ejercicio}{Ejercicio 7. Token, lexema, patrón o atributo \hfill {\small [Esencial]}}
Relaciona cada término con su descripción. Escribe la letra correspondiente.

\begin{center}
\small
\begin{tabular}{|p{0.22\textwidth}|p{0.65\textwidth}|}
\hline
\textbf{Término} & \textbf{Descripción} \tabularnewline
\hline
1. Token & A. Información adicional que distingue la aparición reconocida, por ejemplo el texto de un identificador o el valor de un entero. \tabularnewline
\hline
2. Lexema & B. Descripción de la forma que pueden tener los lexemas asociados con un nombre de token. \tabularnewline
\hline
3. Patrón & C. Secuencia concreta de caracteres reconocida en el programa fuente. \tabularnewline
\hline
4. Atributo & D. Par formado por un nombre de token y un valor de atributo opcional. \tabularnewline
\hline
\end{tabular}
\end{center}

\begin{center}
\begin{tabular}{|c|c|c|c|}
\hline
1 & 2 & 3 & 4 \tabularnewline
\hline
\phantom{D} & \phantom{C} & \phantom{B} & \phantom{A} \tabularnewline[0.45cm]
\hline
\end{tabular}
\end{center}

Completa las relaciones:

\begin{itemize}[leftmargin=*]
    \item El patrón describe una \blank[4cm].
    \item El lexema \blank[4cm] ese patrón.
    \item El token \blank[4cm] el resultado a otra fase.
\end{itemize}
\end{ejercicio}

\ifrespuestas
\begin{respuesta}{Respuesta 7}
1-D, 2-C, 3-B y 4-A.

Relaciones: el patrón describe una \textbf{forma}; el lexema \textbf{satisface} ese patrón; el token \textbf{comunica} el resultado a otra fase.
\end{respuesta}
\fi

\begin{ejercicio}{Ejercicio 8. Analizar un fragmento ilustrativo \hfill {\small [Esencial]}}
Considera el fragmento \texttt{total = 25;}. No supongas que pertenece a \textsc{IMP}. Completa la tabla con los datos del banco.

\begin{center}
\small
\begin{tabular}{|c|c|c|c|}
\hline
\texttt{PYC} & \texttt{ID} & \texttt{ENTERO} & \texttt{ASIG} \tabularnewline
\hline
\(\langle\texttt{ENTERO},25\rangle\) & \(\langle\texttt{PYC}\rangle\) & \(\langle\texttt{ID},\texttt{total}\rangle\) & \(\langle\texttt{ASIG}\rangle\) \tabularnewline
\hline
\end{tabular}
\end{center}

\begin{center}
\small
\renewcommand{\arraystretch}{1.35}
\begin{tabular}{|p{0.15\textwidth}|p{0.17\textwidth}|p{0.34\textwidth}|p{0.21\textwidth}|}
\hline
\textbf{Lexema} & \textbf{Nombre} & \textbf{Patrón informal} & \textbf{Token} \tabularnewline
\hline
\texttt{total} & & Letra seguida opcionalmente de letras o dígitos. & \tabularnewline[0.72cm]
\hline
\texttt{=} & & El carácter \texttt{=}. & \tabularnewline[0.72cm]
\hline
\texttt{25} & & Secuencia de uno o más dígitos. & \tabularnewline[0.72cm]
\hline
\texttt{;} & & El carácter \texttt{;}. & \tabularnewline[0.72cm]
\hline
\end{tabular}
\end{center}

Responde:

\begin{enumerate}[leftmargin=*]
    \item ¿Qué tokens necesitan atributo en este ejemplo y para qué sirve?

    \ifrespuestas\else
    \blank[14cm]

    \blank[14cm]
    \fi

    \item ¿Por qué \texttt{total} es un lexema y no un patrón?

    \ifrespuestas\else
    \blank[14cm]

    \blank[14cm]
    \fi
\end{enumerate}
\end{ejercicio}

\ifrespuestas
\begin{respuesta}{Respuesta 8}
Por filas:

\begin{center}
\small
\renewcommand{\arraystretch}{1.25}
\begin{tabular}{|p{0.15\textwidth}|p{0.17\textwidth}|p{0.34\textwidth}|p{0.21\textwidth}|}
\hline
\textbf{Lexema} & \textbf{Nombre} & \textbf{Patrón informal} & \textbf{Token} \tabularnewline
\hline
\texttt{total} & \texttt{ID} & Letra seguida opcionalmente de letras o dígitos. & \(\langle\texttt{ID},\texttt{total}\rangle\) \tabularnewline
\hline
\texttt{=} & \texttt{ASIG} & El carácter \texttt{=}. & \(\langle\texttt{ASIG}\rangle\) \tabularnewline
\hline
\texttt{25} & \texttt{ENTERO} & Secuencia de uno o más dígitos. & \(\langle\texttt{ENTERO},25\rangle\) \tabularnewline
\hline
\texttt{;} & \texttt{PYC} & El carácter \texttt{;}. & \(\langle\texttt{PYC}\rangle\) \tabularnewline
\hline
\end{tabular}
\end{center}

\begin{enumerate}[leftmargin=*]
    \item \texttt{ID} necesita el atributo \texttt{total} para distinguir ese identificador de otros. \texttt{ENTERO} necesita el atributo 25 para conservar el valor o su representación. En esta descripción, \texttt{ASIG} y \texttt{PYC} no necesitan atributo.
    \item \texttt{total} es la secuencia concreta encontrada en la entrada. El patrón es la descripción general que pueden satisfacer muchos lexemas, por ejemplo ``letra seguida opcionalmente de letras o dígitos''.
\end{enumerate}
\end{respuesta}
\fi

\section*{\cf{verde587246}{5. Comunicación y revisión integradora}}

\begin{ejercicio}{Ejercicio 9. Solicitar el siguiente token \hfill {\small [Esencial]}}
Completa el párrafo con las expresiones del banco. Cada expresión se usa una sola vez.

\begin{center}
\small
\begin{tabular}{|c|c|c|c|c|}
\hline
atributo & lexema & caracteres & token & nombre de token \tabularnewline
\hline
\end{tabular}
\end{center}

Cuando el analizador sintáctico necesita otra unidad, solicita el siguiente \blank[3cm]. El analizador léxico consume los \blank[3cm] necesarios y los agrupa en un \blank[3cm]. Después devuelve un \blank[3.7cm] y, cuando se requiere, un \blank[3cm].

Responde brevemente.

\begin{enumerate}[leftmargin=*]
    \item ¿Es correcto eliminar siempre todos los espacios y saltos de línea? ¿Por qué?

    \ifrespuestas\else
    \blank[14cm]

    \blank[14cm]
    \fi

    \item ¿Qué debe especificar el equipo antes de que Lark pueda producir tokens?

    \ifrespuestas\else
    \blank[14cm]

    \blank[14cm]
    \fi

    \item ¿Cómo puede observarse el analizador léxico de manera independiente con Lark?

    \ifrespuestas\else
    \blank[14cm]

    \blank[14cm]
    \fi
\end{enumerate}
\end{ejercicio}

\ifrespuestas
\begin{respuesta}{Respuesta 9}
En orden: token; caracteres; lexema; nombre de token; atributo.

\begin{enumerate}[leftmargin=*]
    \item No. Deben aplicarse las reglas del lenguaje; en algunos lenguajes los espacios o saltos de línea pueden ser significativos.
    \item Debe especificar las categorías que se reconocerán y el patrón o regla que describe los lexemas de cada una.
    \item Puede utilizarse Lark con \texttt{parser=None}, \texttt{lexer="basic"} y el método \texttt{lex()} para recorrer una entrada y observar los tokens producidos.
\end{enumerate}
\end{respuesta}
\fi

\ifrespuestas\else\newpage\fi
\begin{ejercicio}{Ejercicio 10. Detectar y corregir confusiones \hfill {\small [Extensión]}}
Cada afirmación contiene un error. Subraya la parte incorrecta y reescribe la oración de forma adecuada.

\begin{enumerate}[leftmargin=*]
    \item Un lexema y un patrón son lo mismo: ambos son la secuencia concreta encontrada en el programa.

    \ifrespuestas\else
    \blank[14cm]

    \blank[14cm]
    \fi

    \item El analizador léxico recibe tokens y produce caracteres para el analizador sintáctico.

    \ifrespuestas\else
    \blank[14cm]

    \blank[14cm]
    \fi

    \item Si dos lexemas pertenecen al token \texttt{ID}, entonces deben contener exactamente los mismos caracteres.

    \ifrespuestas\else
    \blank[14cm]

    \blank[14cm]
    \fi

    \item La optimización puede reparar un programa sintácticamente inválido antes de generar código.

    \ifrespuestas\else
    \blank[14cm]

    \blank[14cm]
    \fi

    \item El programa \textsc{C} generado por el compilador del curso ya es directamente un ejecutable de la máquina.

    \ifrespuestas\else
    \blank[14cm]

    \blank[14cm]
    \fi

    \item Lark decide automáticamente qué categorías necesita el lenguaje, aunque el equipo no las especifique.

    \ifrespuestas\else
    \blank[14cm]

    \blank[14cm]
    \fi
\end{enumerate}
\end{ejercicio}

\ifrespuestas
\begin{respuesta}{Respuesta 10}
Correcciones esperadas:
\begin{enumerate}[leftmargin=*]
    \item Un \textbf{lexema} es la secuencia concreta encontrada; un \textbf{patrón} describe la forma de los lexemas de un nombre de token.
    \item El analizador léxico recibe \textbf{caracteres} y produce una secuencia de \textbf{tokens} para el analizador sintáctico.
    \item Dos lexemas pueden corresponder a \texttt{ID} porque satisfacen el mismo patrón, aunque tengan caracteres y atributos distintos.
    \item La optimización recibe una representación válida y la transforma sin cambiar el comportamiento definido; no repara la sintaxis del programa fuente.
    \item El programa \textsc{C} es el destino de nuestro compilador, pero necesita un compilador de \textsc{C} para continuar la cadena hacia un ejecutable.
    \item El equipo debe especificar las categorías y reglas léxicas; Lark aplica esa especificación para producir tokens.
\end{enumerate}
\end{respuesta}
\fi

\end{document}
